\clearpage
\section*{Do the source rules measure the original construction?}
\textit{September 15, 2026. Agent-drafted first-stage population review requested by Jacob.}

The existing rules account for the numerical values in all 13,677
retained construction observations, but they do not consistently
establish that those values describe the original construction. Of the
13,183 observations without recorded overrides, 12,383 reproduce all
four fields from an identified assessment or the declared aggregation of
its components; 636 substitute permit application year for an Assessor
year one year earlier; 162 recover older area measurements; and two fill
missing development-record apartment counts from permit evidence. These
categories partition the ordinary observations. All 494 retained
correction records reproduce their explicit numerical inputs. This is
an accounting result, not independent confirmation of each field's
accuracy or coverage.

The multiple-card rule has a concrete coverage gap. All 339 ordinary
multiple-card observations reproduce their selected cards' unit and
floor sums and take the repeated parcel land once. However, 40 have
additional positive-area cards in the same assessment outside the
selected construction episode, including 13 main-sample observations.
Most extra cards have years before the study period. They could
represent older buildings or obsolete records. The numerator counts
selected construction while the denominator uses parcel land; a
complete selected episode need not cover the whole parcel. These cases
have not been adjudicated or added to the existing flag list.

Among 614 ordinary commercial observations, 452 take all four fields
from the selected assessment. Another 158 take earlier reported areas,
and four combine earlier component records. Twenty of these 162 have a
different count in the earlier area source. Their arithmetic follows
the code; matching parcel membership alone does not establish that old
land and current building measurements describe the same original
construction. The known Madison and Waveland errors also show that
using all fields from one current assessment does not solve the
original-construction question.

The development route retains 64 ordinary observations: 62 use a
positive apartment field and two fill from a permit count. The
Assessor's dictionary defines that apartment field for small multifamily
classes 211/212, so positivity on a development record does not establish
a completed-building count. The 37 completed-condo observations without
recorded overrides instead reproduce their residential-unit membership
and agreeing building fields. All 321 ordinary tied-parcel observations
reproduce their selected counts, floor areas and parcel-land sums.
These internal agreements do not independently establish physical lot
coverage.

The application-year substitution already affects 636 observations,
including 201 in the main sample. Application year is a proxy rather
than a recorded completion date; this review has not established that
those 636 dates are wrong. A separate provenance-label defect affects
five ordinary records: floor recovery identifies a 2025 source row
while the output retains a 2022 assessment-year label. Their other
values agree under the recovery rule, so this finding alone does not
establish wrong density measurements.

The full-date permit check finds permit records sharing the parcel base
(the first ten digits) for 7,295
projects and a full residential new-building candidate for 4,002.
There are 505 with a single completed candidate issued after the
assigned year and 62 with later residential conversion or count-change
text. These are chronology leads, not 505 wrong years or 62 inflated
counts. A parcel can contain another building or a later construction
episode, and some later conversions are already correctly excluded.
No permit date or count was substituted by this audit.

\textit{Scope and reproduction.} Production values, the 218 existing
flags, estimates and manuscript are unchanged. The existing flags are
reconciled below; no sample of unflagged buildings has been drawn or
reviewed. The audit producer
\path{tasks/audits/construction_measurement_review/code/check_source_rules.R}
writes project-level and source-route tables with SaveData reports. It
uses the preserved ordinary history, completed-condo history, correction
file and full permit history. Make holds the existing production and
screening outputs fixed for this first-stage review. The HTML report
and \path{source_rules.md} explain the distinctions between accounting,
source coverage and substantive confirmation. The Assessor's field
scope is documented at
\url{https://github.com/ccao-data/model-res-avm#features-used}.

\textit{Reconcile the existing flags before new research.} Reading all
64 recorded decisions for their actual subject separates earlier
evidence from mere source preferences. Of the original 218 flags, 47
are addressed by earlier documented decisions, ten have current-count
support from the recent investigation, and six properties already lie
outside the common FAR and DUPAC sample. Eight have confirmed count
errors, with corrections not applied. Sixteen have recently reviewed
but open measurement questions, seven have old decisions that leave
the flagged question open, 72 have source disagreements still to
reconcile, and 52 have repetition or ratio clues without a case
adjudication. These categories partition the original 218; their
main-sample counts are respectively 12, 9, 0, 5, 9, 2, 6 and 17,
totaling 60.

The earlier decision for Belmont's two 27-home buildings cites
separate permit revisions, so equal counts do not by themselves
reopen that decision. The mixed-use floor area at 440 N Halsted also
has an explicit explanation. Conversely, expired conflicting permits
do not prove newer Assessor counts correct, and Optima Signature's
count decision does not settle its repeated land area. Eight of the
64 older decisions rely on a source preference without independently
settling counts; one, Lawrence, is already included among the recently
reviewed open cases. Reusing a decision addresses its specific flag,
not every field of the building. This is an evidence reconciliation,
not completion of all remaining case adjudications. The renderer
produces \path{flag_reconciliation.csv} with a data report, using the
screen, recent findings and \path{prior_decision_review.csv}. Production
data remain byte-identical to the preserved analysis snapshot.

\clearpage
\section*{Review of all 87 unchanged-area count disagreements}
\textit{September 15, 2026. Agent recommendations following the requested case review; not adopted corrections.}

All 87 retained projects with a positive unchanged-area count-conflict
indicator were reviewed using preserved assessment histories, permit
records, earlier decisions and available external building evidence.
They include 27 observations in the main 500-foot sample. The count and
land recommendations are to retain 45 current choices, correct 24
counts (one also needs different land), remove six conversions,
renovations or additions, retire one obsolete duplicate, and leave
eleven unresolved or provisional. The main-sample counts in those
categories are respectively 15, 6, 4, zero and two. These categories
partition the reviewed records, but do not certify all their variables:
twenty of the same records have construction-timing questions and three
may carry forward an old building's floor area. The evidence table
states those concerns alongside the count and land recommendations.

There is no defensible unconditional oldest-assessment or
newest-assessment rule. Original full-building permits and completed
building accounts often corroborate an Assessor count, while revisions
and later conversion permits explain other changes. At 5439 N Lincoln,
the evidence distinguishes two dwellings from commercial space; at
4351 N Lincoln it supports three dwellings rather than either reported
four or five. This does not justify automatically subtracting one from
all mixed-use apartment fields. At 1214 S Sawyer a 2014 permit explicitly
converts the original two-flat to a single-family house. At 1664 W
Division the original permits specify 60 apartments and a later permit
converts a fitness center into another apartment. The construction
observation should describe the original building rather than its
later count.

At Madison, independent building evidence supports nine apartments at
1100 and eighteen at 1048, retaining their reported lots of 4,599 and
10,000 square feet. Eighty-one describes the community, and even the
older 24 at 1048 is unsupported. The contemporary account reports 2013
completion, which also needs reconciliation with the selected 2012.
Wolf Point West's developer reports 509 apartments and a 40,653-square-foot
site instead of the repeated 178,133-square-foot development total;
its initial occupancy is reported as 2016 rather than the selected
2014. Wolf Point East's individual lot remains unresolved. Circa 922
has 104 newly built apartments and 45 existing lofts, so its combined
site cannot automatically be paired with the new-building count.

The six removal recommendations concern 2656 W North (warehouse
conversion), 2425 N Rockwell, 1939 N Whipple, 2223 N Bissell, 116 N
Elizabeth and 3440 N Southport (additions or renovations). Wholly new
replacement buildings remain eligible in principle, but identical
floor numbers across demolition and replacement do not verify the
new building's floor area. That issue is retained explicitly for 2742 W
38th, 1428 N Cleveland and 2045 W Ohio. At 660 E 42nd, consecutive old
and successor assessment histories share an address; the successor's
preserved 2009 parcel point falls inside the old 2006 polygon, 6.15 feet
from its centroid in CRS 3435. This supports retiring the obsolete
observation, while the successor's original count and year remain
unresolved. No land was calculated from those polygons.

The eleven unresolved or provisional observations are Lawrence, Lake,
Circa 922, Wolf Point East, Michigan, Crystal, the two Sawyer houses,
the successor at 660 E 42nd, 552 W 43rd and 2935--2937 W 63rd. Only
Lawrence and Michigan are currently in the main sample. The two Sawyer
counts of two remain reasonable Assessor-based preferences but lack
independent original-building confirmation. Crystal and 43rd have
evidence of pre-period construction despite later selected years. The
63rd Street permit may describe only part of the combined assessment;
neither its count nor a replacement denominator is proposed without
settling that coverage.

\textit{Proposed rule and implementation.} Flag unchanged-area count
changes; establish the original completed building; separate dwellings
from shops and later conversions; and retain source-reported land for
the same building or coherent group. An issued or expired permit alone
does not establish completion. The interpretation needed here should
be documented in the existing correction input once approved, applied
once in ordinary cleaning. This review introduces no production task,
no address-specific correction code and no automatic numerical choice
between contradictory snapshots.

\textit{Research record and verification.}
\path{tasks/audits/construction_measurement_review/count_land_review.csv}
contains all 87 recommendations, current and proposed measurements,
assessment row IDs and histories, exact relevant permit descriptions,
source URLs, limitations and an unadopted indicator. The existing
renderer checks unique keys, exact coverage of the population indicator,
agreement with current count/land/sample membership, and unadopted
status. It writes the HTML review and a standard data report. The
report was rebuilt through the audit Makefile with the six existing
screen/source-check outputs held fixed. Production data, recorded
production decisions, estimates and manuscript remain unchanged; no
unflagged-building sample was drawn. Sources and the report-only build
command are recorded in \path{count_land_review.md}.

\clearpage
\section*{Follow-up recommendations for the same 87 records}
\textit{September 15, 2026. Further source review; all recommendations remain unadopted.}

The follow-up covers thirty of the same 87 records: the union of the
previously unresolved or provisional cases and the overlapping date
and floor-area questions. The proposed treatment across all 87 is now
to keep 65 after specified corrections, retain two Sawyer observations
with an explicit count qualification, exclude six additions,
renovations or conversions, exclude one partial ownership record and
one lodging building, place two buildings before the study period,
retire one obsolete duplicate, and withhold nine records whose needed
measurements or construction years remain unresolved. These categories
sum to 87. Among the proposed retained observations, 24 counts,
14 construction years and one lot area would change; the changes
overlap. Current main-sample membership comprises 19 proposed retained
observations, four non-new-construction exclusions, one partial-building
exclusion, one lodging exclusion and two unresolved records. Corrected
years could change historical boundary assignment, so these are not
post-adoption sample counts.

The two Madison buildings have documented completion in 2013, with
nine and eighteen apartments. Other completed-building evidence supports
2011 for 544 W Oak, 2015 for Eight O Five, 2017 for Aur\'elien, 2016 for
Centrum Wicker Park and Luxe on Chicago, 2016 for Wolf Point West,
2008 for Burnham Pointe, and 2017 for 6614 W 65th. A June 2021 listing
offers newly completed apartments at 2910 S Archer; the exact-PIN permit
identifies that address rather than the selected 2904. Earlier assessment,
retail-buildout and rental evidence support three apartments in 2021
at 3542 N Southport. Six apartments and a 2014 completion proxy are
supported at 3408 N Milwaukee. Van Buren's 2014 remains an earlier
Assessor proxy amid inconsistent marketing, rather than a located
occupancy certificate. Superior's rental history establishes completed
use by mid-2022, supporting retention of the Assessor's 2022 year and
three dwellings despite an expired permit status.

The follow-up also shows why initial permit counts cannot always settle
a dispute. Aur\'elien's developer page has completed-project data of
368 apartments in July 2017 but unrevised proposal prose stating 373.
The full permit and completed-project data support retaining 368,
while the conflicting older assessment and publicity are preserved in
the case record. At 1428 N Cleveland, the single-family permit concealed
a use question: an October 2016 visitor describes a newly built
bed-and-breakfast, The Sono, with seven guest rooms and owners living
on site. The operator and seller descriptions corroborate lodging use.
The recommendation is to exclude it from ordinary residential
construction, not count its guest rooms as dwellings. This remains a
sample-definition judgment requiring agreement.

The nine withholding recommendations concern Lawrence, Lake,
Circa 922, Wolf Point East, Jackson, 38th, Ohio, the successor at 42nd,
and the 63rd Street parcel pair. Additional sources did not establish
the missing completion year, building coverage, floor measure or
individual lot. At Ohio, seller marketing reports an estimated area
including finished lower-level rooms; it does not independently replace
the Assessor-comparable floor area. At 38th, the located old listing
describes the demolished building. Withholding FAR removes these two
from the agreed common FAR-and-DUPAC sample even though one dwelling
and reported land remain supported. No map-calculated or allocated
land area is proposed. Crystal's original building is dated 2005 in
the early assessment, and the 43rd Street house already appears in
2005 with year built 2004; both should remain outside the study period.
The Michigan record groups only seven holdings within a much larger
condominium tower and is not a whole-building observation.

The two Sawyer properties are explicitly qualified retention judgments.
Each has an initial count of three and then two in all subsequent
preserved assessments, without a located original-building permit or
independent seller account. Two is the recommended Assessor-based
choice, not independently established truth or a general majority-vote
rule. The research pass is complete for the bounded follow-up, but
adoption, conservative withholding and the lodging judgment still need
agreement. This does not adjudicate all other population flags or begin
the deferred unflagged-building sample.

\textit{Reproduction and production boundary.} Findings, URLs and
proposed values were added to the existing 87-row audit evidence table.
The existing renderer reports proposed sample treatment and displays
the follow-up evidence, while preserving source histories and the
unadopted indicator. The audit report and this logbook were rebuilt
through their Makefiles against the preserved upstream outputs.
Production correction and analysis files remain unchanged. Eventual
approved decisions belong in the existing committed correction CSV,
joined once by ordinary cleaning; this follow-up adds no production
scripts or tasks.


\clearpage
\section*{Adoption of the apartment-count review}
\textit{September 15, 2026. Approved decisions entered in the existing correction table.}

Jacob approved incorporating the recommendations for the 87 reviewed records.
The adopted treatment retains 67, including the two qualified Sawyer counts,
and leaves twenty outside the common FAR-and-DUPAC sample. Six describe
additions, renovations or conversions rather than new buildings; one is a
partial ownership record, one is lodging, two predate the study period and
one is an obsolete duplicate. Nine lack a supported measurement or completion
year. For the two replacement houses with unverified floor area, the dwelling
count and reported land remain usable for DUPAC, but FAR is withheld; the
common sample therefore excludes them. The Sawyer records retain two apartments
each as documented Assessor-based judgments, not independently verified counts.

The adoption check corrected one mistake in the earlier review itself.
The proposed 112 apartments and 2011 for 531 W Division came from the separate
544 W Oak building. The Oak permit carried the Division parcel number, but
its street address and coordinate identify the other Parkside phase. The
completed original permit 100186134 at 545 W Division specifies 111 apartments;
the Chicago Housing Authority's May 2025 Cabrini NOW completion table reports
111 units and 2009 for Parkside IB-Rental. The adopted Division observation
therefore uses 111 and retains its selected Assessor year 2009, floor area,
lot and location. This supersedes the 112/2011 recommendation above. Among
the retained 67 records, 25 dwelling counts, thirteen years and one land area
change relative to the pre-adoption data; these changes overlap.

Seven other retained buildings required explicit location references after
count or year changes. LaSalle, Chicago Avenue, Fulton Market and Milwaukee
use completed original-building permits. Wolf Point West, Archer and Southport
use the preserved 2025 points for their exact County parcel numbers after
checking the building identity. These are parcel-location proxies, not
building-footprint measurements. Wolf Point West's address and recorded owner
identify the West building; its original permit point lies farther east on
the shared site. The developer's reported individual lot, rather than any
area calculated from these coordinates, supplies its density denominator.
Archer's County property-address field still says 2904, while the exact-parcel
permit and taxpayer mailing address identify 2910. These source choices and
limitations are stated with the corrections.

The existing production correction CSV receives fifty new rows and two
revisions; the other 35 reviewed observations need no redundant override.
All 740 unrelated earlier correction rows are unchanged. Ordinary cleaning
applies this table once; no production script or Makefile rule changed.
The Make rebuild using preserved inputs yields 13,664 construction records,
compared with 13,677 before adoption. Eligibility for both density measures
falls from 13,635 to 13,615. All 67 retained review records have the specified
measurements and usable locations; the other twenty are outside the common
sample. Every unreviewed project and boundary row is unchanged. A second
Make check requires no computation. The audit and logbook were rebuilt;
estimation inputs, estimates and manuscript remain unchanged. Audit snapshots
preserve the original discrepancies. This completes adoption of the bounded
87-record review, not adjudication of other flags or the deferred unflagged
sample.

\clearpage
\section*{Density estimates after the measurement corrections}
\textit{September 15, 2026. Rerun requested after committing the reviewed decisions.}

The corrections weaken the multifamily DUPAC estimate materially, while the
two all-construction estimates remain negative and significant at five percent.
Holding the original stringency scores and regression specifications fixed,
the estimates comparing the first 100 feet on each side of the boundary are:

\begin{center}
\begin{tabular}{lrrr}
\hline
Outcome & Before (\%) & Corrected (\%) & Corrected $p$ \\
\hline
All construction: FAR & $-9.24$ & $-9.82$ & .025 \\
All construction: DUPAC & $-13.70$ & $-12.82$ & .048 \\
Multifamily: FAR & $-12.85$ & $-11.51$ & .069 \\
Multifamily: DUPAC & $-17.36$ & $-11.68$ & .175 \\
\hline
\end{tabular}
\end{center}

Effects are $100[\exp(\hat\beta)-1]$. The fitted all-construction sample falls
from 4,023 to 4,016 projects. Multifamily remains at 863, with changed membership.
Within the eligible 500-foot data, nine observations exit, two enter following
location corrections, and two retained houses become multifamily after their
original two-unit counts are restored. Fixed-effect omissions account for the
difference between these eligible counts and fitted counts. All changes in
the rebuilt analysis data belong to the reviewed properties; every unreviewed
row remains unchanged. The citywide common-density sample has 13,615 projects,
including 2,570 multifamily projects.

At 1100 and 1048 W Madison, the adopted counts are nine and eighteen apartments,
with 2013 completion. DUPAC is now 85.24 and 78.41, using the unchanged reported
lots. Both remain assigned to Burnett opposite Fioretti. These corrections are
part of the full rerun; the table above does not isolate their individual
contributions to the change.

Ordering the corrected observations by raw average log processing time gives
all-construction FAR and DUPAC estimates of $-2.24\%$ and $-0.70\%$, and
multifamily estimates of $-3.76\%$ and $-5.22\%$. None is significant at ten
percent. Dropping disagreements instead gives $-10.36\%$, $-11.55\%$,
$-8.00\%$ and $-12.48\%$; only all-construction FAR reaches ten percent
($p=.085$). Thus the reviewed measurement errors explain part of the original
multifamily difference, but the sensitivity to the ordering remains. This
exercise does not establish which score measures alderman behavior better.
All uncertainty statements condition on the estimated ordering and selected
sample; score-estimation uncertainty is not included.

The existing appendix checks were also rebuilt. All eight placebo coefficients
remain statistically indistinguishable from zero at ten percent. All eight
donut point estimates remain negative, although their significance varies.
Removing each project's matched permits from the existing score calculation
leaves the density estimates almost unchanged: all-construction FAR moves from
$-.103$ to $-.104$ log points, and the other coefficients agree to three decimal
places. The existing score-gap restrictions were rerun as well. This is the
implemented project-specific permit exclusion check, not a new full first-stage
leave-out estimator or combined bootstrap.

\textit{Reproduction.} The reviewed decisions were committed and pushed as
\texttt{7b0cf59c}. Make rebuilt the two analysis-data producers and the existing
main density, density appendix and density score-robustness tasks using the
preserved scores, controls, zoning sources, permit history and boundary inputs.
This is a downstream rerun on the existing checkout, not a fresh-clone build.
The raw-log audit now reads the corrected production analysis data and matches
its alderman pairs to the preserved citywide score/time table. Its prior density
estimates and figure are saved in its \texttt{records/} directory; the older
driver audit continues to use those pre-correction inputs. Sales, rent and permit
outcome estimates are retained unchanged. No production R script or Makefile
was added or modified. The working-paper text and compiled PDF remain unchanged.

\clearpage
\begin{center}
\includegraphics[width=\textwidth]{../tasks/audits/raw_log_score_sensitivity/output/raw_log_effects.pdf}
\end{center}
